\section{Weighted scaling and the quadratic model}
\label{sec:scaling}

We now turn the boundary data into a global model of $U$.  There are two
distinct convergence statements to retain: compact subsets of the proposed
limit eventually lie in the scaled domains, and points forbidden by the
limit are approached by excluded points.  The latter will prevent the
limiting biholomorphism from acquiring an unwanted image component or a
boundary value.
Boundary scaling is a method developed by
Pinchuk~\cite[Section~2, pp.~73--75]{Pinchuk1981};
Frankel's compact-quotient argument likewise studies limits of normalized
automorphisms~\cite[Section~3]{Frankel1989}. We give the set convergence
and the forward and inverse compactness arguments in the form required
by the possibly higher-codimensional boundary cell chosen here.

Fix the point $p$, cell $S$, old coordinates $(z_0,w_0)$, excluded cone $F$,
and $k$ support functions $h_1,\ldots,h_k$ supplied by
Proposition~\ref{prop:supported-point}.  We retain the old parameters
$s=(z_0,\operatorname{Re}w_0)$ and offsets
$v=\operatorname{Im}w_0-\phi(s)$ when invoking
Proposition~\ref{prop:boundary-data}.  By
Proposition~\ref{prop:noncollapse}, $F\ne\R^k$.

Choose translated, complex-linearly adjusted coordinates $(z,w)$ centered
at $p$ such that
\[
 T_pS=\{\operatorname{Im}w=0\},
 \qquad (z,w)\in\C^m\times\C^k,\qquad m=n-k.
\]
Let $\Phi$ denote the inverse coordinate chart, with $\Phi(0)=p$, and
shrink its neighborhood $\Omega$ in $V$ when necessary.  Put
\[
 \delta_t(Z,W)=(\sqrt t\,Z,tW),\qquad
 D_t=\delta_t^{-1}\bigl(\Phi^{-1}(U\cap\Omega)\bigr)
 \quad(t>0).
\]
Every fixed compact subset of $\C^n$ lies in the domain of
$\Phi\circ\delta_t$ for all sufficiently small $t$; membership in $D_t$
then means exactly that its image lies in $U$.

\begin{proposition}[Two-sided convergence of the scaled sets]
\label{prop:scaling-limit}
There are an invertible real linear map $L:\R^k\to\R^k$ and a vector
$Q$ of real quadratic forms on $\C^m$ such that, with
\[
 R(Z,W)=L(\operatorname{Im}W)-Q(Z),\qquad
 \mathcal O=\{(Z,W):R(Z,W)\notin F\},
\]
the following hold.
\begin{enumerate}[label=\textup{(\roman*)}]
\item Every compact subset of $\mathcal O$ is contained in $D_t$ for all
sufficiently small $t>0$.
\item For every $\xi\in\C^n$ with $R(\xi)\in F$ and every sequence
$t_j\downarrow0$, there are $\xi_j\notin D_{t_j}$ such that
$\xi_j\to\xi$.  These excluded points can be chosen within the scaled
coordinate charts.
\end{enumerate}
\end{proposition}

\begin{proof}
Regard the old offset $v$ as a real analytic function of the new
coordinates.  Its differential has kernel $T_pS$, so
$dv_0(z,w)=L(\operatorname{Im}w)$ with $L$ invertible.  Taylor expansion
therefore gives
\begin{equation}
\label{eq:residual-scaling}
 R_t(Z,W):=t^{-1}v\bigl(\Phi(\sqrt t\,Z,tW)\bigr)
 \longrightarrow L(\operatorname{Im}W)-Q(Z)
\end{equation}
uniformly on compact sets.  Here $-Q$ is the part of the quadratic Taylor
term involving only $z$; the mixed and normal quadratic terms disappear
after division by $t$.  The old parameter
$s_t(Z,W)=s(\Phi(\sqrt t\,Z,tW))$ tends uniformly on compact sets to
$s_*$.

For clarity, the distance convergence in
Proposition~\ref{prop:boundary-data} says that
\[
 d(s,t;u)=\min\{1,\operatorname{dist}
       (u,\R^k\setminus t^{-1}A_s)\}
 \longrightarrow d_*(u)
 \quad\text{as }(s,t)\longrightarrow(s_*,0^+),
\]
locally uniformly in $u$, with $F=\{d_*=0\}$.  In particular it applies
to moving vectors $u_j\to u$, as well as to moving parameters $s_j$.
If assertion~\textup{(i)} failed, a compactness argument would give
$\xi_j\to\xi\in\mathcal O$ and $t_j\downarrow0$ with
$\xi_j\notin D_{t_j}$.  The corresponding offsets
$R_{t_j}(\xi_j)$ are excluded from $t_j^{-1}A_{s_{t_j}(\xi_j)}$.
Their distances to the excluded sets are thus zero, whereas their limit
has distance value $d_*(R(\xi))>0$.  This is a contradiction.

For assertion~\textup{(ii)}, set $q_j=\Phi(\delta_{t_j}\xi)$ and retain
its old parameter $s_j=s(q_j)$.  Since $R_{t_j}(\xi)\to R(\xi)\in F$,
the same joint convergence supplies
\[
 u_j\in\R^k\setminus t_j^{-1}A_{s_j},
 \qquad u_j-R_{t_j}(\xi)\longrightarrow0.
\]
In the old chart change only the imaginary normal coordinate, replacing
the offset $v(q_j)$ by $t_ju_j$.  The resulting point $q'_j$ is outside
$U$, has the same $s_j$, and differs from $q_j$ by $o(t_j)$ in old
coordinates.  It remains in $\Omega$ for large $j$.  The fixed coordinate
change is smooth, so the change is also $o(t_j)$ in the new coordinates.
Consequently
\[
 \xi_j:=\delta_{t_j}^{-1}\Phi^{-1}(q'_j)
 \longrightarrow\xi,
 \qquad \xi_j\notin D_{t_j}.
\]
Indeed an $o(t_j)$ change becomes $o(\sqrt{t_j})$ in the $Z$ coordinates
and $o(1)$ in the $W$ coordinates.  All points used here approach $p$,
so the local cutoffs in the definition of $A_s$ do not affect the argument.
\end{proof}

Choose a connected component $C$ of the nonempty open cone
$L^{-1}(\R^k\setminus F)$, and set
\begin{equation}
\label{eq:preliminary-model}
 P_2=L^{-1}Q,\qquad
 D_0=\{(Z,W):\operatorname{Im}W-P_2(Z)\in C\}.
\end{equation}
Every positive dilation preserves $C$: for $a\in C$, the path
$r\mapsto ra$ stays in the complement cone and connects $a$ to any
positive multiple.  The real coordinate change
\[
 (Z,W)\longmapsto
       (Z,\operatorname{Re}W,\operatorname{Im}W-P_2(Z))
\]
identifies $D_0$ with $\C^m\times\R^k\times C$.  Thus $D_0$ is connected
and is precisely one connected component of $\mathcal O$.

\begin{theorem}[Quadratic model]
\label{thm:quadratic-model}
Under the hypotheses of Theorem~\ref{thm:main}, with $n\ge1$, there are
integers $1\le k\le n$, $m=n-k$, a nonempty connected open cone
$C\subset\R^k$, and a vector $H$ of Hermitian forms on $\C^m$ such that
$U$ is biholomorphic to
\begin{equation}
\label{eq:model}
 D=\{(z,w)\in\C^m\times\C^k:
               \operatorname{Im}w-H(z,z)\in C\}.
\end{equation}
Here $C$ is invariant under every positive dilation, and $H$ is complex
linear in its first argument.  There are independent real linear
functionals $\ell_1,\ldots,\ell_k$ on $\R^k$ and a constant $c>0$ such
that
\begin{equation}
\label{eq:strict-supports}
 \ell_a(y)>0\quad(y\in C),\qquad
 \ell_a(H(z,z))\ge c\lVert z\rVert^2
 \quad(z\in\C^m),\qquad 1\le a\le k.
\end{equation}
The original action transfers to a properly discontinuous action on $D$
with compact quotient.  No freeness assumption is required.
\end{theorem}

\begin{proof}
We first prove $U\simeq D_0$, and then remove the pure quadratic terms
and pass the support inequalities to the limit.

\paragraph{Normalization and forward bounds.}
Choose $a_0\in C$ and write $b_0=(0,ia_0)\in D_0$.  By
Proposition~\ref{prop:scaling-limit}, for any sufficiently small sequence
$t_j\downarrow0$ the points
$p_j=\Phi(0,it_ja_0)$ belong to $U$.  The compact covering set $K$ gives
\[
 \gamma_j(k_j)=p_j,\qquad \gamma_j\in\Gamma,\quad k_j\in K.
\]
After extraction, $k_j\to k_*\in K$.

Each support $h_a$ maps $U$ into the unit disc, satisfies $h_a(p)=1$, and
has $h_a(p_j)=1+O(t_j)$.  For a fixed compact subset $E\subset U$,
$d_U(q,k_j)$ is bounded uniformly for $q\in E$.  Schwarz contraction
for the chain-of-discs distance implies that
$h_a(\gamma_jq)$ is at uniformly bounded hyperbolic distance from
$h_a(p_j)$.  If two disc points $\alpha,\beta$ have pseudohyperbolic
distance at most $\rho<1$, then
\[
 |\beta-\alpha|
 \le \frac{\rho(1-|\alpha|^2)}{1-\rho|\alpha|}
 \le \frac{2\rho}{1-\rho}(1-|\alpha|).
\]
It follows that, uniformly on $E$,
\begin{equation}
\label{eq:scaled-support-bound}
 h_a(\gamma_jq)=1+O_E(t_j),\qquad 1\le a\le k.
\end{equation}
The peak inequality in Equation~\eqref{eq:peak} now gives
\[
 \lVert\gamma_jq-p\rVert^2
 \le -\log|h_a(\gamma_jq)|=O_E(t_j).
\]
Here and below the norm of points in $V$ is the norm in its fixed affine
embedding.  In particular, the images of $E$ eventually lie in $\Omega$,
and their new coordinates satisfy $(z,w)=O_E(\sqrt{t_j})$.

The $k$ supports sharpen the bound in all the normal directions.
Their real logarithmic differentials annihilate
$T_pS=\{\operatorname{Im}w=0\}$ and are independent.  Therefore, taking
the local logarithms normalized by $\log h_a(p)=0$, we have
\begin{equation}
\label{eq:normal-linear-supports}
 d(\log h_a)_p=i\ell_a(dw),
\end{equation}
where $\ell_1,\ldots,\ell_k$ are independent real rows, extended
complex linearly when applied to $w$.  Indeed a complex linear form whose
real part vanishes on all $z$ and real $w$ has no $z$ term and purely
imaginary $w$ coefficients.  Taylor expansion, the preliminary
$O_E(\sqrt{t_j})$ bound, and
Equation~\eqref{eq:scaled-support-bound} give
\[
 i\ell_a(w)=O_E(t_j),\qquad 1\le a\le k.
\]
Invertibility of the matrix of these rows gives $w=O_E(t_j)$.
Consequently the holomorphic maps
\[
 f_j=\delta_{t_j}^{-1}\circ\Phi^{-1}\circ\gamma_j
\]
are eventually defined on a neighborhood of every compact subset of $U$
and are locally uniformly bounded.  A diagonal normal-family extraction
gives a holomorphic limit $f:U\to\C^n$.

\paragraph{Inverse compactness.}
For each $j$ the inverse coordinate map
\[
 g_j=\gamma_j^{-1}\circ\Phi\circ\delta_{t_j}:D_{t_j}\longrightarrow U
\]
is defined on all of $D_{t_j}$ and satisfies $g_j(b_0)=k_j$.  Every compact
subset of $D_0$ is eventually contained in its domain, by
Proposition~\ref{prop:scaling-limit}.  More importantly, for every compact
$E\subset D_0$ there is a compact $K_E\subset U$ containing $g_j(E)$
for all sufficiently large $j$.

To see this last assertion, cover $E$ by finitely many balls
$B(x_\nu,r_\nu)$ with $\overline B(x_\nu,2r_\nu)\subset D_0$.
Connectedness and small balls along paths give a fixed finite disc chain
from $b_0$ to each $x_\nu$, with all closed disc images in $D_0$.
For $x\in B(x_\nu,r_\nu)$, append the disc
\[
 \zeta\longmapsto x_\nu+2\zeta(x-x_\nu),
\]
whose parameters $0$ and $1/2$ join $x_\nu$ to $x$.  All these chains
lie in one compact subset of $D_0$, and their total Poincar\'e lengths
have a common bound.  That compact subset is contained in $D_{t_j}$
for all large $j$, so composition with $g_j$ proves
\[
 d_U(g_j(x),k_j)\le M_E\qquad(x\in E)
\]
with $M_E$ independent of $j$.  Since $k_j\in K$ and $d_U$ is proper by
Proposition~\ref{prop:proper-distance}, this is the claimed compactness.

Montel's theorem in the bounded ambient affine coordinates, followed by
diagonal extraction, now gives a holomorphic map $g:D_0\to U$ with
$g_j\to g$ locally uniformly.  The compact containment just proved puts
the limit in $U$, rather than merely in its closure.  It also allows the
composition identities to pass to the limit: on each compact subset of
$D_0$, the points $g_j(x)$ lie in a common compact subset of $U$, where
$f_j\to f$.  Thus
\begin{equation}
\label{eq:one-composition}
 f\circ g=\operatorname{id}_{D_0}.
\end{equation}
In particular, $f$ has full rank at every point of $g(D_0)$.

This identity has not yet located the rest of $f(U)$ or proved global
injectivity. The next step uses the excluded-point conclusion of
Proposition~\ref{prop:scaling-limit} to settle both issues.

\paragraph{Excluding additional image points.}
The maps $f_j$ have nonzero Jacobian wherever they are defined.  Their
local Jacobian determinants converge to that of $f$.  The zero-free
limit principle says that a limit of zero-free holomorphic functions is
either zero-free or identically zero; in several variables this follows
by restriction to complex lines.  Thus the full-rank locus and its
complement are both open: a zero of a local limiting determinant forces
that determinant to vanish throughout its coordinate ball.
Equation~\eqref{eq:one-composition} makes the full-rank locus nonempty,
so connectedness of $U$ implies that $f$ has nonzero Jacobian everywhere.

Fix $q\in U$.  Local inverse stability gives a neighborhood $B$ of $f(q)$
such that
\begin{equation}
\label{eq:stable-image-ball}
 B\subset f_j(V_q)\subset D_{t_j}
\end{equation}
for all sufficiently large $j$, where $V_q\Subset U$ is a fixed
coordinate neighborhood of $q$.  Here is an explicit justification of
the stability assertion.  Center source coordinates at $q$ and normalize
the target by the inverse of $df_q$.  On a sufficiently small closed
source ball the normalized derivative of $f$ differs from the identity
by less than $1/4$.  Local uniform convergence of the holomorphic maps
gives the same bound with $1/2$ for $f_j$ on that ball.  For every target
point in a fixed smaller ball about $f(q)$, the equation $f_j(x)=y$ is
then solved by a contraction of the closed source ball into itself,
since $f_j(q)\to f(q)$.  This proves
Equation~\eqref{eq:stable-image-ball}.

If $R(f(q))\in F$, assertion~\textup{(ii)} of
Proposition~\ref{prop:scaling-limit} supplies points outside $D_{t_j}$
converging to $f(q)$, contradicting
Equation~\eqref{eq:stable-image-ball}.  Thus $f(U)\subset\mathcal O$.
The set $f(U)$ is connected and contains $D_0$ by
Equation~\eqref{eq:one-composition}, so it lies in that same connected
component: $f(U)=D_0$.  Finally, for each $q\in U$ the points $f_j(q)$
eventually lie in a compact neighborhood of $f(q)$ inside $D_0$.
Passing to the limit in $g_j(f_j(q))=q$ yields $g(f(q))=q$.  We have proved
that $f:U\to D_0$ and $g:D_0\to U$ are inverse biholomorphisms.

\paragraph{The Hermitian form and its supports.}
The real quadratic vector $P_2$ has a unique decomposition
\[
 P_2(z)=H(z,z)+\operatorname{Re}B(z),
\]
where $H$ is a vector of Hermitian forms and $B$ a vector of holomorphic
homogeneous quadratic polynomials.  The triangular holomorphic change
\[
 (Z,W)\longmapsto(z,w)=(Z,W-iB(Z))
\]
is globally invertible and transforms $D_0$ into $D$ of
Equation~\eqref{eq:model}.

The functionals in Equation~\eqref{eq:normal-linear-supports} retain
their independence.  We claim that throughout $D$ they satisfy
\begin{equation}
\label{eq:model-support-bound}
 \ell_a(\operatorname{Im}w)\ge c\lVert z\rVert^2,
 \qquad 1\le a\le k,
\end{equation}
for a common $c>0$.  To prove this, fix $(z,w)\in D$ and consider the
unscaled point
\[
 q_t=\Phi\bigl(\sqrt t\,z,
                 t(w+iB(z))\bigr).
\]
It lies in $U$ for every sufficiently small $t$, by
Proposition~\ref{prop:scaling-limit} applied before the triangular
change.  The local holomorphic logarithm of a support has expansion
\[
 \log h_a(\Phi(z,w))
   =i\ell_a(w)+A_a(z)
      +O\bigl(\lVert z\rVert^3+
               \lVert z\rVert\lVert w\rVert+
               \lVert w\rVert^2\bigr),
\]
where $A_a$ is holomorphic homogeneous quadratic.  Dividing the peak
inequality $-\log|h_a(q_t)|\ge\lVert q_t-p\rVert^2$ by $t$ and passing
to the limit gives
\begin{equation}
\label{eq:unaveraged-support}
 \ell_a(\operatorname{Im}w)
       +\operatorname{Re}\bigl(\ell_a(B(z))-A_a(z)\bigr)
 \ge \lVert Jz\rVert^2,
 \qquad J=d_z\Phi_0.
\end{equation}
The linear map $J$ is injective.  Thus, when $m>0$, its right-hand side
is bounded below by $c\lVert z\rVert^2$ with $c>0$.  When $m=0$, the
same inequality holds with any positive $c$, since $z=0$.

For each real $\theta$, the point $(e^{i\theta}z,w)$ also lies in $D$,
because $H(e^{i\theta}z,e^{i\theta}z)=H(z,z)$.  Average
Equation~\eqref{eq:unaveraged-support} over $\theta$.  The holomorphic
quadratic terms have average zero, whereas
$\lVert J(e^{i\theta}z)\rVert=\lVert Jz\rVert$.  This proves
Equation~\eqref{eq:model-support-bound}.

Putting $z=0$ gives $\ell_a(y)\ge0$ on $C$.  Since $C$ is open and
$\ell_a\ne0$, equality at a point of $C$ would produce a nearby negative
value.  Hence $\ell_a(y)>0$ for every $y\in C$.  Next fix $y\in C$.
For every $r>0$,
\[
 (z,i(H(z,z)+ry))\in D.
\]
Apply Equation~\eqref{eq:model-support-bound} and let $r\downarrow0$ to
obtain $\ell_a(H(z,z))\ge c\lVert z\rVert^2$.  This proves both parts of
Equation~\eqref{eq:strict-supports}.  Conjugating the original group by
the resulting biholomorphism transfers proper discontinuity and
cocompactness to $D$, completing the proof.
\end{proof}
